% =====================================================================
%  Lecture 2026-09-24

%  CORRECTED SOURCE ERRORS
%   [p1] algebraic basis defined in a normed space  ->  in a vector
%        space; the norm plays no role, as the source itself warns on p2
%   [p1] partial sum s_n := sum_{n=1}^N x_n  ->  s_N := sum_{n=1}^N x_n
%   [p1] Schauder basis given as a set  ->  a sequence, since the
%        series depends on the order
%   [p1] ||x - s_N||_p = sum_{i>N} |a_i|^p  ->  ||x - s_N||_p^p
%   [p2] complex Fourier system 1, e^{ix}, e^{2ix}, ... with the sum
%        from k = 1  ->  e^{ikx} for all k in Z, symmetric partial sums
%   [p2] ||f - sum c_k e^{ikx}||_2 = integral of |.|^2  ->  its square
%   [p2] rational coefficients for separability  ->  Q + iQ when K = C
%   [p4] finite-dimensional example: ||Tx|| bounded by ||sum a_i e_i||,
%        without T, and the bound |a_i| <= ||x|| of the boxed revision
%        ->  ||sum a_i T e_i||, with the coefficient bound from Lemma 3.2
%   [p5] f_x(y) = <x,y>  ->  f_x(y) = sum x_i y_i, linear over R and C
%   [p5] sup [f(x) g(x)] without absolute values  ->  sup |f(x) g(x)|
%   [p6] evaluation functional called I in two places  ->  E
%   [p6] integral kernel example: X = C([0,1] x [0,1])  ->  the operator
%        acts on C[0,1]; only the kernel lives on the square
%   [p8] converse: ||x|| <= delta => ||Tx|| <= 1 from continuity at 0,
%        which only gives it for ||x|| < delta  ->  delta/2 used, bound
%        2/delta
%   [p10] B(X,X) a Banach algebra for normed X  ->  for Banach X
%
%  ADDITIONS
%   [p1] E not an algebraic basis of ell_p: the witness (2^{-n})
%   [p1] E a Schauder basis of ell_p: uniqueness via the coordinate
%        bound |y_j| <= ||y||_p
%   [p1] ell_infty: any representation of (1,1,...) must have all
%        coefficients 1, so none converges
%   [p1] the definition of series convergence moved before the Schauder
%        basis, which uses it
%   [p2] ell_infty not separable: distinct 0-1 sequences give distinct
%        points of the dense set
%   [p2] Schauder basis implies separable: the two-step density estimate;
%        promoted from remark to proposition, since a corollary rests on it
%   [p2, p3] Zorn's lemma stated before the theorem that uses it, with
%        chain and maximal element glossed; in the proof: the family is
%        non-empty, the union of a chain is an independent upper bound,
%        a maximal independent set spans
%   [p3] absolute convergence: the index N from convergence of
%        sum ||x_i||
%   [p4] ||Tx|| <= ||T|| ||x||: the case x = 0
%   [p5] the restated examples I and II merged into the first two
%   [p5] multiplication on ell_p: termwise estimate, which also shows
%        M_x y in ell_p; p range 1 <= p < infinity
%   [p7] integral operator: T_K f continuous, by uniform continuity of K
%   [p7] bounded implies continuous: the case T = 0
%   [p8] corollary on continuity at 0: the converse used only that
%   [p8] B(X,Y) is a vector space: supplied in place of the homework box
%   [p8] norm axioms: ||T|| = 0 => T = 0 via ||Tx|| <= ||T|| ||x||;
%        homogeneity written out
%   [p9] completeness of B(X,Y): linearity of the limit, and boundedness
%        of T from T = T_n - (T_n - T)
%   [p10] Banach algebra glossed in the sentence that uses it
%
%  REPORTED BUT UNMODIFIED GAPS
%   [p2] completeness of the trigonometric systems in L_2 stated without
%        proof; so is the failure of pointwise convergence
%   [p7] incompleteness of the polynomials under ||.||_infty asserted
%
%  OMITTED
%   [p5, p6, p7] homework boxes ||M_f|| = ||f||_infty,
%        ||M_x|| = ||x||_infty, ||I|| = 1, ||T_K|| = ||K||_infty, as
%        unsolved exercises. The last is false: ||T_K|| =
%        sup_x int |K(x,y)| dy, and K(x,y) = y has ||K||_infty = 1,
%        ||T_K|| = 1/2.
%   [p6] the matrix sketches, including "T^T = T", and the kernel square;
%        the matrix analogy is kept as a formula
%   [p10] the invertibility sketch
%   [p4, p7, p9] the writer's asides and crossed-out lines
% =====================================================================

% [p1]
\section{Bases}

\begin{definition}[Algebraic basis]\label{def:hamel}
Let $X$ be a vector space. A set $S \subset X$ is an \emph{algebraic basis},
or \emph{Hamel basis}, if $S$ is linearly independent and every $x \in X$ has
a representation $x = \sum_{i \in I} \alpha_i s_i$ with $s_i \in S$,
$\alpha_i \in \K$ and $I$ finite.
\end{definition}

\begin{example}\label{ex:E-not-hamel}
Let $1 \le p < \infty$ and $e_n = (0, \dots, 0, 1, 0, \dots)$, with the $1$
in position $n$. The set $E = \set{e_1, e_2, \dots}$ is not an algebraic
basis of $\lp{p}$: a finite combination of the $e_n$ has only finitely many
non-zero coordinates, while $(2^{-n})_n \in \lp{p}$ has infinitely many.
\end{example}

\begin{definition}[Convergent series]\label{def:series}
In a normed space $X$ a series $\sum_{n=1}^{\infty} x_n$ \emph{converges} to
$x \in X$ if its partial sums $s_N := \sum_{n=1}^{N} x_n$ satisfy
$s_N \to x$.
\end{definition}

\begin{definition}[Schauder basis]\label{def:schauder}
In a Banach space $X$ a sequence $(s_n)$ is a \emph{Schauder basis} if every
$x \in X$ has a unique representation $x = \sum_{i=1}^{\infty} \alpha_i s_i$
with $\alpha_i \in \K$, the series converging in the sense of
Definition~\ref{def:series}.
\end{definition}

\begin{example}\label{ex:E-schauder}
$E = (e_n)$ is a Schauder basis of $\lp{p}$ for $1 \le p < \infty$.
\end{example}

\begin{proof}
Let $x = (\alpha_1, \alpha_2, \dots) \in \lp{p}$ and
$s_N = \sum_{i=1}^{N} \alpha_i e_i = (\alpha_1, \dots, \alpha_N, 0, \dots)$.
Then
\[
  \norm{x - s_N}_p^p
  = \norm{(0, \dots, 0, \alpha_{N+1}, \alpha_{N+2}, \dots)}_p^p
  = \sum_{i=N+1}^{\infty} \abs{\alpha_i}^p \longrightarrow 0,
\]
being the tail of a convergent series, so $x = \sum_i \alpha_i e_i$. For
uniqueness, suppose $x = \sum_i \beta_i e_i$. The $j$-th coordinate of the
$N$-th partial sum is $\beta_j$ once $N \ge j$, and a coordinate satisfies
$\abs{y_j} \le \norm{y}_p$, so coordinates pass to the limit and
$\beta_j = \alpha_j$.
\end{proof}

\begin{remark}\label{rem:E-linfty}
$E$ is not a Schauder basis of $\lp{\infty}$. Take $x = (1, 1, \dots)$.
Since $\abs{y_j} \le \norm{y}_\infty$, the argument above shows that any
representation $x = \sum_i \beta_i e_i$ must have $\beta_i = 1$ for all $i$.
But then
\[
  \norm{x - s_N}_\infty = \norm{(0, \dots, 0, 1, 1, \dots)}_\infty = 1
\]
for every $N$, so the partial sums do not converge to $x$, and $x$ has no
representation at all.
\end{remark}

% [p2]
\begin{remark}\label{rem:linfty-nonsep}
$\lp{\infty}$ is not separable. The sequences with every term $0$ or $1$ form
an uncountable set, and any two distinct ones have $\norm{\cdot}_\infty$
distance $1$. The balls of radius $\frac{1}{3}$ around them are therefore
pairwise disjoint: a common point would put the two centres at distance at
most $\frac{2}{3}$. A dense set meets every one of these balls, and meets
distinct balls in distinct points, so it is uncountable.
\end{remark}

\begin{nfigure}{Around the uncountably many $0$--$1$ sequences, the balls of
radius $\frac{1}{3}$ in $\lp{\infty}$ are pairwise disjoint, and a dense set
needs a point in each.}
\begin{tikzpicture}[scale=0.9]
  \draw[rulegrey] (0,0) ellipse (2.4 and 1.25);
  \foreach \x/\y in {-1.6/0.2, -0.9/-0.6, -0.6/0.55, 0.1/-0.2,
                      0.6/0.6, 1.0/-0.7, 1.6/0.15} {
    \draw[accent] (\x,\y) circle (0.2);
    \fill[accent] (\x,\y) circle (0.03);
  }
  \node[rulegrey] at (2.75,1.05) {$\lp{\infty}$};
\end{tikzpicture}
\end{nfigure}

\begin{example}\label{ex:fourier}
The Fourier systems are Schauder bases. In the real space $\Lp{2}[0,2\pi]$ the
system $1, \cos x, \sin x, \cos 2x, \sin 2x, \dots$ is one: every $f$ has a
unique representation
\[
  f = a_0 + \sum_{k=1}^{\infty} \bigl(a_k \cos(kx) + b_k \sin(kx)\bigr).
\]
In the complex space $\Lp{2}[0,2\pi]$ the system $e^{ikx}$, $k \in \Z$,
enumerated as $1, e^{ix}, e^{-ix}, e^{2ix}, e^{-2ix}, \dots$, is one: every
$f$ has a unique representation $f = \sum_{k=-\infty}^{\infty} c_k e^{ikx}$,
meaning
\[
  \Bigl\lVert f - \sum_{k=-N}^{N} c_k e^{ikx} \Bigr\rVert_2^2
  = \int_0^{2\pi} \Bigl\lvert f(x) - \sum_{k=-N}^{N} c_k e^{ikx}
    \Bigr\rvert^2 \dd x
  \xrightarrow[N \to \infty]{} 0 .
\]
Convergence in $\norm{\cdot}_2$ does not guarantee pointwise convergence.
\end{example}

\begin{proposition}\label{prop:schauder-sep}
A Banach space with a Schauder basis is separable.
\end{proposition}

\begin{proof}
Let $(s_i)$ be a Schauder basis and let $D$ be the set of finite sums
$\sum_{i=1}^{N} q_i s_i$ with $N \in \N$ and $q_i \in \Q$, or
$q_i \in \Q + i\Q$ when $\K = \C$. For each $N$ these sums form a countable
set, so $D$ is countable. Given $x = \sum_i \alpha_i s_i$ and $\eps > 0$,
choose $N$ with $\norm{x - \sum_{i=1}^{N} \alpha_i s_i} < \frac{\eps}{2}$, and
then rational $q_i$ with
$\abs{\alpha_i - q_i}\,\norm{s_i} < \frac{\eps}{2N}$ for $i \le N$. The
triangle inequality gives $\norm{x - \sum_{i=1}^{N} q_i s_i} < \eps$, so $D$
is dense.
\end{proof}

\begin{corollary}\label{cor:linfty-no-schauder}
$\lp{\infty}$ has no Schauder basis.
\end{corollary}

\begin{proof}
It is not separable, by Remark~\ref{rem:linfty-nonsep}, so
Proposition~\ref{prop:schauder-sep} excludes a Schauder basis.
\end{proof}

\begin{lemma}[Zorn]\label{lem:zorn}
Let $(P, \le)$ be a partially ordered set in which every chain, that is every
totally ordered subset, has an upper bound in $P$. Then $P$ has a maximal
element, an $m \in P$ such that $m \le p$ implies $p = m$.
\end{lemma}

% [p3]
\begin{remark}
The \emph{axiom of choice} states that $\prod_{i \in I} A_i$ is non-empty
whenever every $A_i$ is non-empty, that is, that there is a family
$(a_i)_{i \in I}$ with $a_i \in A_i$ for every $i$.
\end{remark}

% [p2]
\begin{theorem}\label{thm:hamel-exists}
Every vector space has an algebraic basis.
\end{theorem}

\begin{remark}
An algebraic basis has no connection to any norm on the space.
\end{remark}

\begin{proof}[Proof of Theorem~\ref{thm:hamel-exists}]
Let $P$ be the family of linearly independent subsets of the vector space
$X$, and write $I \le J$ if $I \subseteq J$. This is a partial order: it is
reflexive, $I \le I$; antisymmetric, $I \le J$ and $J \le I$ give $I = J$;
and transitive, $I \le J$ and $J \le K$ give $I \le K$. The family is
non-empty, since $\emptyset \in P$.

% [p3]
Let $\mathcal{C} \subset P$ be a chain. Its union $U$ is linearly
independent: finitely many $u_1, \dots, u_k \in U$ lie in members
$C_1, \dots, C_k$ of $\mathcal{C}$, which are totally ordered, so the largest
of them contains all $u_j$, and these are therefore independent. Thus
$U \in P$ is an upper bound of $\mathcal{C}$; for the empty chain,
$\emptyset$ is one. By Lemma~\ref{lem:zorn} there is a maximal independent
set $S$.

$S$ spans $X$. Otherwise there is an $x \notin \operatorname{span} S$, and
$S \cup \set{x}$ is independent: in a vanishing combination the coefficient
of $x$ must be $0$, since otherwise $x \in \operatorname{span} S$, and then
the remaining coefficients vanish by independence of $S$. This contradicts
maximality, so $S$ is an algebraic basis.
\end{proof}

\begin{definition}[Absolute convergence]\label{def:abs-conv}
A series $\sum_{i=1}^{\infty} x_i$ in a normed space $X$ is
\emph{absolutely convergent} if $\sum_{i=1}^{\infty} \norm{x_i} < \infty$.
\end{definition}

\begin{lemma}[Absolute convergence]\label{lem:abs-conv}
If $X$ is a Banach space and $\sum_{i=1}^{\infty} x_i$ is absolutely
convergent, then $\sum_{i=1}^{\infty} x_i$ converges.
\end{lemma}

\begin{proof}
As $X$ is complete, it suffices to show that the partial sums
$s_N = \sum_{i=1}^{N} x_i$ form a Cauchy sequence. Let $\eps > 0$. The real
series $\sum_i \norm{x_i}$ converges, so there is an $N$ with
$\sum_{i=m+1}^{n} \norm{x_i} < \eps$ for all $n > m > N$. For such $n, m$,
\[
  \norm{s_n - s_m} = \Bigl\lVert \sum_{i=m+1}^{n} x_i \Bigr\rVert
  \le \sum_{i=m+1}^{n} \norm{x_i} < \eps . \qedhere
\]
\end{proof}

% [p4]
\section{Bounded linear operators}

\begin{definition}[Linear operator]\label{def:linear-op}
Let $X$ and $Y$ be vector spaces over $\K$. A map $T \colon X \to Y$ is
\emph{linear} if $T(x_1 + x_2) = Tx_1 + Tx_2$ for all $x_1, x_2 \in X$ and
$T(\alpha x) = \alpha\, Tx$ for all $\alpha \in \K$ and $x \in X$.
\end{definition}

\begin{definition}[Bounded operator, operator norm]\label{def:bounded-op}
A linear operator $T \colon (X, \norm{\cdot}) \to (Y, \norm{\cdot})$ between
normed spaces is \emph{bounded} if
$\sup_{\norm{x} \le 1} \norm{Tx} < \infty$. In that case
\[
  \norm{T} = \sup_{\norm{x} \le 1} \norm{Tx}
\]
is the \emph{operator norm} of $T$.
\end{definition}

\begin{nfigure}{A bounded operator maps the closed unit ball to a bounded set,
here an ellipse; $\norm{T}$ is the radius of the smallest ball about $0$
containing the image.}
\begin{tikzpicture}[scale=0.85]
  \draw[rulegrey, -{Stealth[length=1.6mm]}] (-1.6,0) -- (1.6,0);
  \draw[rulegrey, -{Stealth[length=1.6mm]}] (0,-1.4) -- (0,1.4);
  \draw[accent, thick] (0,0) circle (0.8);
  \node[accent] at (1.25,-0.95) {$B(0,1)$};
  \node at (-1.3,1.3) {$X$};
  \draw[-{Stealth[length=2mm]}] (1.8,0.9) to[bend left=25]
    node[above] {$T$} (3.4,0.9);
  \begin{scope}[xshift=5.4cm]
    \draw[rulegrey, -{Stealth[length=1.6mm]}] (-1.9,0) -- (1.9,0);
    \draw[rulegrey, -{Stealth[length=1.6mm]}] (0,-1.6) -- (0,1.6);
    \draw[thmaccent, thick, rotate=30] (0,0) ellipse (1.4 and 0.6);
    \node[thmaccent] at (1.75,1.2) {$T(B(0,1))$};
    \node at (-1.6,1.5) {$Y$};
  \end{scope}
\end{tikzpicture}
\end{nfigure}

\begin{remark}\label{rem:op-norm-bound}
For every $x \in X$, $\norm{Tx} \le \norm{T}\,\norm{x}$.
\end{remark}

\begin{proof}
For $x = 0$ both sides vanish, since $T0 = 0$ by linearity. For $x \ne 0$,
linearity gives $Tx = \norm{x}\, T\bigl(x / \norm{x}\bigr)$, and
$x / \norm{x}$ has norm $1$, so
\[
  \norm{Tx} = \norm{x}\, \Bigl\lVert T\Bigl(\frac{x}{\norm{x}}\Bigr)
  \Bigr\rVert
  \le \norm{x} \sup_{\norm{x'} \le 1} \norm{Tx'} = \norm{x}\,\norm{T} .
  \qedhere
\]
\end{proof}

\begin{example}\label{ex:finite-dim-bounded}
Every linear operator $T \colon X \to Y$ on a finite-dimensional normed space
$X$ is bounded. In particular every linear $T \colon \K^n \to \K^m$, which is
given by an $m \times n$ matrix, is bounded.
\end{example}

\begin{proof}
Let $e_1, \dots, e_n$ be a basis of $X$. By Lemma~\ref{lem:independent}
there is a $C > 0$ with $\norm{\sum_i \alpha_i e_i} \ge C \sum_i
\abs{\alpha_i}$ for all coefficients. For $x = \sum_{i=1}^{n} \alpha_i e_i$,
\[
  \norm{Tx} = \Bigl\lVert \sum_{i=1}^{n} \alpha_i\, T e_i \Bigr\rVert
  \le \sum_{i=1}^{n} \abs{\alpha_i}\,\norm{T e_i}
  \le \Bigl(\max_{1 \le i \le n} \norm{T e_i}\Bigr)
      \sum_{i=1}^{n} \abs{\alpha_i}
  \le \frac{1}{C} \max_{1 \le i \le n} \norm{T e_i}\; \norm{x},
\]
so $\norm{T} \le C^{-1} \max_i \norm{T e_i} < \infty$.
\end{proof}

% [p5]
\begin{example}\label{ex:fx}
For fixed $x \in \K^n$ let $f_x \colon \K^n \to \K$,
$f_x(y) = \sum_{i=1}^{n} x_i y_i$. Then $f_x$ is linear, and bounded by
Example~\ref{ex:finite-dim-bounded} with $m = 1$.
\end{example}

\begin{definition}[Linear functional]\label{def:functional}
A linear operator from $X$ to $\K$ is a \emph{linear functional}.
\end{definition}

\begin{example}[Multiplication operator on $C{[0,1]}$]\label{ex:mult-C}
Fix $f \in \Cont{0,1}$ and let $M_f \colon \Cont{0,1} \to \Cont{0,1}$,
$g \mapsto f g$, with the norm $\norm{\cdot}_\infty$. Then $M_f$ is linear
and bounded, with $\norm{M_f} \le \norm{f}_\infty$.
\end{example}

\begin{proof}
$\norm{f}_\infty < \infty$ because $f$ is continuous on the compact $[0,1]$.
For $\norm{g}_\infty \le 1$,
\[
  \norm{M_f g}_\infty = \sup_{x \in [0,1]} \abs{f(x) g(x)}
  \le \norm{f}_\infty \norm{g}_\infty \le \norm{f}_\infty . \qedhere
\]
\end{proof}

\begin{example}[Multiplication operator on $\ell_p$]\label{ex:mult-lp}
Let $1 \le p < \infty$, fix $x = (x_1, x_2, \dots) \in \lp{\infty}$, and let
$M_x \colon \lp{p} \to \lp{p}$,
$y = (y_1, y_2, \dots) \mapsto (x_1 y_1, x_2 y_2, \dots)$. Then $M_x$ is
linear and bounded, with $\norm{M_x} \le \norm{x}_\infty$.
\end{example}

\begin{proof}
Linearity is immediate. Termwise
$\abs{x_i y_i}^p \le \norm{x}_\infty^p \abs{y_i}^p$, and summing gives
$\norm{M_x y}_p \le \norm{x}_\infty \norm{y}_p$. This shows first that
$M_x y \in \lp{p}$, so that $M_x$ is well defined, and then that
$\norm{M_x y}_p \le \norm{x}_\infty$ whenever $\norm{y}_p \le 1$.
\end{proof}

\begin{remark}
Multiplication operators are the simplest operators. They are the
counterpart of diagonal matrices: $D = \operatorname{diag}(\alpha_1, \dots,
\alpha_n)$ acts by $y \mapsto Dy = (\alpha_1 y_1, \dots, \alpha_n y_n)$,
which is multiplication by $(\alpha_1, \dots, \alpha_n)$, and diagonalizing a
matrix in linear algebra reduces it to such an operator.
\end{remark}

% [p6]
\begin{example}[Integral functional]\label{ex:integral-functional}
Let $I \colon \Cont{0,1} \to \K$, $I(f) = \int_0^1 f(t) \dd t$, with the
norm $\norm{\cdot}_\infty$. Then $I$ is linear and bounded, with
$\norm{I} \le 1$: for $\norm{f}_\infty \le 1$,
\[
  \abs{I f} = \Bigl\lvert \int_0^1 f(t) \dd t \Bigr\rvert
  \le \int_0^1 \abs{f(t)} \dd t \le \int_0^1 \dd t = 1 .
\]
\end{example}

\begin{example}[Evaluation functional]\label{ex:evaluation}
Fix $t_0 \in [0,1]$ and let $E \colon \Cont{0,1} \to \K$, $E f = f(t_0)$,
with the norm $\norm{\cdot}_\infty$. Then $E$ is linear and bounded, with
$\norm{E} \le 1$: for $\norm{f}_\infty \le 1$,
$\abs{E f} = \abs{f(t_0)} \le \norm{f}_\infty \le 1$.
\end{example}

\begin{example}[Integral operator]\label{ex:integral-kernel}
Fix a \emph{kernel} $K \in C([0,1] \times [0,1])$ and let
$T_K \colon \Cont{0,1} \to \Cont{0,1}$,
\[
  (T_K f)(x) = \int_0^1 K(x,y) f(y) \dd y ,
\]
with the norm $\norm{\cdot}_\infty$ on both sides. This is the continuous
analogue of matrix multiplication, $(Kf)_j = \sum_{m=1}^{n} K_{jm} f_m$.
Then $T_K$ is linear and bounded, with $\norm{T_K} \le \norm{K}_\infty$,
where $\norm{K}_\infty = \max_{[0,1]^2} \abs{K}$ is finite because $K$ is
continuous on a compact set.
\end{example}

% [p7]
\begin{proof}
$T_K f$ is continuous: $K$ is uniformly continuous on the compact square, and
$\abs{(T_K f)(x) - (T_K f)(x')}
\le \sup_{y} \abs{K(x,y) - K(x',y)}\,\norm{f}_\infty$. Linearity is
immediate. For $\norm{f}_\infty \le 1$,
\[
  \norm{T_K f}_\infty
  = \sup_{x \in [0,1]} \Bigl\lvert \int_0^1 K(x,y) f(y) \dd y \Bigr\rvert
  \le \sup_{x \in [0,1]} \int_0^1 \abs{K(x,y)}\,\abs{f(y)} \dd y
  \le \norm{K}_\infty \int_0^1 \dd y = \norm{K}_\infty . \qedhere
\]
\end{proof}

\begin{example}[Differentiation]\label{ex:differentiation}
Let $X$ be the space of polynomials on $[0,1]$ with the norm
$\norm{\cdot}_\infty$, a subspace of $\Cont{0,1}$; it is not a Banach space.
Let $D \colon X \to X$, $f \mapsto f'$. Then $D$ is linear but not bounded:
for $f_n(t) = t^n$ one has $\norm{f_n}_\infty = 1$ for all $n$, while
$\norm{D f_n}_\infty = \norm{n t^{n-1}}_\infty = n \to \infty$.
\end{example}

\begin{remark}
As a rule of thumb, integration is bounded and differentiation is not.
\end{remark}

\section{The space of bounded operators}

\begin{theorem}\label{thm:bounded-continuous}
A linear operator $T \colon (X, \norm{\cdot}) \to (Y, \norm{\cdot})$ is
bounded if and only if it is continuous.
\end{theorem}

\begin{proof}
Assume $T$ is bounded. If $T = 0$ it is continuous, so let $\norm{T} > 0$.
Given $\eps > 0$ put $\delta = \eps / \norm{T}$. For all $x_1, x_2 \in X$
with $\norm{x_1 - x_2} < \delta$, Remark~\ref{rem:op-norm-bound} gives
\[
  \norm{T x_1 - T x_2} = \norm{T(x_1 - x_2)}
  \le \norm{T}\,\norm{x_1 - x_2}
  < \norm{T} \frac{\eps}{\norm{T}} = \eps .
\]

% [p8]
Assume $T$ is continuous; only continuity at $0$ will be used. Since
$T0 = 0$, the choice $\eps = 1$ gives a $\delta > 0$ with $\norm{Tx} < 1$
whenever $\norm{x} < \delta$. For $\norm{z} \le 1$ the vector
$\frac{\delta}{2} z$ has norm less than $\delta$, so
\[
  \norm{T z} = \frac{2}{\delta} \Bigl\lVert T\Bigl(\frac{\delta}{2} z\Bigr)
  \Bigr\rVert < \frac{2}{\delta} .
\]
Hence $T$ is bounded, with $\norm{T} \le 2/\delta$.
\end{proof}

\begin{corollary}\label{cor:cont-at-0}
A linear operator $T \colon (X, \norm{\cdot}) \to (Y, \norm{\cdot})$ is
continuous everywhere if and only if it is continuous at $0$.
\end{corollary}

\begin{proof}
The second half of the proof of Theorem~\ref{thm:bounded-continuous} used
only continuity at $0$ to obtain boundedness, and the first half turns
boundedness into continuity everywhere.
\end{proof}

\begin{definition}\label{def:BXY}
$\Bop{X}{Y}$ denotes the set of all bounded linear operators from $X$ to $Y$.
\end{definition}

\begin{theorem}\label{thm:BXY-normed}
$\Bop{X}{Y}$ is a vector space under pointwise operations, and the operator
norm is a norm on it. Therefore $(\Bop{X}{Y}, \norm{\cdot})$ is a normed
space.
\end{theorem}

\begin{proof}
Pointwise sums and scalar multiples of linear maps are linear, and the zero
operator is bounded. For $S, T \in \Bop{X}{Y}$, $\alpha \in \K$ and
$\norm{x} \le 1$,
\[
  \norm{(S + T)x} \le \norm{Sx} + \norm{Tx} \le \norm{S} + \norm{T},
  \qquad
  \norm{(\alpha T)x} = \abs{\alpha}\,\norm{Tx} \le \abs{\alpha}\,\norm{T},
\]
so $S + T$ and $\alpha T$ are bounded, and $\Bop{X}{Y}$ is a subspace of
the vector space of linear maps from $X$ to $Y$.

For the norm axioms, $\norm{T} \ge 0$ is clear. If $\norm{T} = 0$, then
Remark~\ref{rem:op-norm-bound} gives $\norm{Tx} \le \norm{T}\,\norm{x} = 0$
for every $x$, so $T = 0$. Homogeneity follows from
$\norm{\alpha T} = \sup_{\norm{x} \le 1} \abs{\alpha}\,\norm{Tx}
= \abs{\alpha}\,\norm{T}$. For the triangle inequality, the triangle
inequality in $Y$ gives
\[
  \norm{T + S} = \sup_{\norm{x} \le 1} \norm{(T + S)x}
  \le \sup_{\norm{x} \le 1} \bigl(\norm{Tx} + \norm{Sx}\bigr)
  \le \sup_{\norm{x} \le 1} \norm{Tx} + \sup_{\norm{x} \le 1} \norm{Sx}
  = \norm{T} + \norm{S} . \qedhere
\]
\end{proof}

% [p9]
\begin{theorem}\label{thm:BXY-banach}
If $(Y, \norm{\cdot})$ is a Banach space, then so is
$(\Bop{X}{Y}, \norm{\cdot})$.
\end{theorem}

\begin{corollary}\label{cor:dual-banach}
$(\Bop{X}{\K}, \norm{\cdot})$ is a Banach space for every normed space $X$.
\end{corollary}

\begin{notation}
$X' = (\Bop{X}{\K}, \norm{\cdot})$ is the \emph{dual space} of $X$.
\end{notation}

\begin{proof}[Proof of Theorem~\ref{thm:BXY-banach}]
Let $(T_n) \subset \Bop{X}{Y}$ be a Cauchy sequence: for every $\eps > 0$
there is an $N$ with $\norm{T_n - T_m} \le \eps$ for all $n, m > N$. For
each $x \in X$, Remark~\ref{rem:op-norm-bound} gives
\[
  \norm{T_n x - T_m x} = \norm{(T_n - T_m)x}
  \le \norm{T_n - T_m}\,\norm{x} \le \eps \norm{x},
\]
and since $\eps$ is arbitrary, $(T_n x)$ is Cauchy in $Y$. As $Y$ is
complete, the pointwise limit $Tx := \lim_{n \to \infty} T_n x$ exists. It
is linear, because
$T(\alpha x + \beta y) = \lim_n (\alpha T_n x + \beta T_n y)
= \alpha Tx + \beta Ty$.

It remains to show that $T$ is bounded and that $T_n \to T$ in norm. Fix
$n > N$ and let $m \to \infty$ in $\norm{T_n x - T_m x} \le \eps\norm{x}$;
continuity of the norm gives
\[
  \norm{(T_n - T)x} \le \eps \norm{x} \qquad \text{for all } x \in X .
\]
Thus $T_n - T$ is bounded with $\norm{T_n - T} \le \eps$, so
$T = T_n - (T_n - T) \in \Bop{X}{Y}$ by Theorem~\ref{thm:BXY-normed}, and
$\norm{T_n - T} \le \eps$ for all $n > N$ means $T_n \to T$.
\end{proof}

\begin{proposition}[Composition]\label{prop:composition}
If $T \in \Bop{Y}{Z}$ and $S \in \Bop{X}{Y}$, then $TS \in \Bop{X}{Z}$ and
$\norm{TS} \le \norm{T}\,\norm{S}$.
\end{proposition}

% [p10]
\begin{proof}
$TS$ is linear as a composition of linear maps. By
Remark~\ref{rem:op-norm-bound},
\[
  \norm{TS} = \sup_{\norm{x} \le 1} \norm{T(Sx)}
  \le \sup_{\norm{x} \le 1} \norm{T}\,\norm{Sx}
  = \norm{T} \sup_{\norm{x} \le 1} \norm{Sx} = \norm{T}\,\norm{S} .
  \qedhere
\]
\end{proof}

\begin{remark}
In particular, for a Banach space $X$ the space $\Bop{X}{X}$ is a Banach
algebra: a Banach space, by Theorem~\ref{thm:BXY-banach}, which is also an
algebra under composition, with $\norm{TS} \le \norm{T}\,\norm{S}$.
\end{remark}

\begin{definition}[Invertible operator]\label{def:invertible}
$T \in \Bop{X}{X}$ is \emph{invertible} if $T$ is a bijection from $X$ to
$X$ and $T^{-1}$ is bounded.
\end{definition}
